\section{Normal Duality and Geometrization}
\label{sec:normal_duality}

\paragraph{Question.}
The fixed-\(\beta\) dual map \(g\mapsto g^\vee=g_\star^2/g\) reflects the scalar coordinate
\(y=\log\chi\) but keeps \(\beta\) unchanged. A natural geometric refinement is to ask whether one
can define a dual representative \((\beta^\perp,g^\perp)\) that simultaneously moves in both
coordinates, with connecting segment orthogonal to the crossover curve.

\paragraph{Crossover geometry in \((\beta,g)\).}
For fixed \(\theta\), define
\begin{equation}
  y(\beta,g,\theta)=\log\chi(\beta,g,\theta),
  \qquad
  \mathcal{C}_\theta:=\{(\beta,g):y=0\}.
\end{equation}
In one-channel scaling,
\begin{equation}
  y(\beta,g,\theta)=2\log g+\log\chi_0(\beta,\theta),
\end{equation}
so
\begin{equation}
  \nabla y
  =\left(\partial_\beta\log\chi_0,\frac{2}{g}\right).
\end{equation}
Hence the Euclidean normal direction to \(\mathcal{C}_\theta\) is generated by \(\nabla y\), and
the tangent slope is
\begin{equation}
  g_\star'(\beta,\theta)=-\frac{g_\star(\beta,\theta)}{2}\,\partial_\beta\log\chi_0(\beta,\theta).
\end{equation}

\begin{definition}[Normal-ray dual representative]
\label{def:normal_dual}
Assume \(y\in C^2\) and \(\nabla y\neq 0\) on \(\mathcal{C}_\theta\). Let \(\mathcal{U}_\theta\) be a
tubular neighborhood where each \(p\in\mathcal{U}_\theta\) has unique closest point
\(\pi(p)\in\mathcal{C}_\theta\), and write
\begin{equation}
  p=\pi(p)+\lambda\,\hat{\nu}(\pi(p)),
\end{equation}
with \(\hat{\nu}\) the unit normal field.
Define \(p^\perp=\mathcal{D}_\perp(p)\) as the point on the same normal ray
\begin{equation}
  p^\perp=\pi(p)+\lambda^\perp \hat{\nu}(\pi(p)),
\end{equation}
such that
\begin{equation}
  y(p^\perp)=-y(p).
\end{equation}
\end{definition}

\begin{theorem}[Local existence and involution of normal duality]
\label{thm:normal_duality}
Under Definition~\ref{def:normal_dual}, there exists \(\delta>0\) such that for
\(|\lambda|<\delta\):
\begin{enumerate}
  \item the scalar \(\lambda^\perp\) exists and is unique;
  \item \(\mathcal{D}_\perp\) is an involution, \(\mathcal{D}_\perp^2=\mathrm{id}\);
  \item the segment joining \(p\) and \(p^\perp\) is orthogonal to \(\mathcal{C}_\theta\) at
  \(\pi(p)\).
\end{enumerate}
\end{theorem}
\begin{proof}
Fix \(c\in\mathcal{C}_\theta\) and define \(\Phi_c(\lambda):=y(c+\lambda\hat{\nu}(c))\).
Since \(y(c)=0\) and \(\Phi_c'(0)=\nabla y(c)\!\cdot\!\hat{\nu}(c)=\|\nabla y(c)\|>0\), \(\Phi_c\)
is strictly monotone on a small interval by continuity. Therefore, for any \(\lambda\) in that
interval, the equation \(\Phi_c(\lambda^\perp)=-\Phi_c(\lambda)\) has unique solution
\(\lambda^\perp\). This defines \(\mathcal{D}_\perp\) locally. Applying the same construction twice
gives \(\Phi_c(\lambda^{\perp\perp})=-\Phi_c(\lambda^\perp)=\Phi_c(\lambda)\), and injectivity of
\(\Phi_c\) yields \(\lambda^{\perp\perp}=\lambda\), hence involution. Orthogonality is immediate
because both points lie on one normal ray at \(c=\pi(p)\).
\end{proof}

\begin{corollary}[Simultaneous \texorpdfstring{\((\beta,g)\)}{(beta,g)} dual representative]
\label{cor:simultaneous_beta_g_dual}
The map \(\mathcal{D}_\perp\) defines a local simultaneous dual
\begin{equation}
  (\beta,g)\mapsto(\beta^\perp,g^\perp),
\end{equation}
with
\begin{equation}
  y(\beta^\perp,g^\perp,\theta)=-y(\beta,g,\theta).
\end{equation}
Generically both coordinates change. In the special case where
\(\partial_\beta\log\chi_0\approx 0\) at the crossover point, the normal is nearly vertical and the
map reduces to the fixed-\(\beta\) duality to leading order.
\end{corollary}

\paragraph{Interpretation.}
This provides a geometrization of the strong-weak exchange: instead of viewing duality only as an
algebraic inversion in \(g\) at fixed \(\beta\), one may view it as local reflection across the
crossover manifold along its normal bundle. The orientation-gauge \(Z_2\) then acts as an
independent discrete fiber operation on top of this base-manifold reflection.

